\FloatBarrier
\newpage
\section{反常积分*}
\label{sec:calculus-improper-integrals}

此前积分的区间有限，被积函数也有界。若范围延伸到无穷远，或者函数在端部无界，就先在正常的有限范围内积分，再研究这些积分的极限。无穷累积是否有限，由这个新的极限过程决定。

\subsection{无穷区间上的积分}

\begin{definition}
设\(f\)在\([a,+\infty)\)的每个有限闭子区间上可积。若有限极限存在，定义
\begin{equation*}
 \int_a^{+\infty}f(x)\,dx=\lim_{R\to+\infty}\int_a^Rf(x)\,dx,
\end{equation*}
并称该{\heiti{}反常积分收敛}；否则称其发散。左侧无穷区间的定义类似。
\end{definition}

两端都无穷时，选取一个有限分点\(c\)，分别定义\(\int_{-\infty}^cf\)与\(\int_c^{+\infty}f\)；只有两者都收敛，才定义其和为\(\int_{-\infty}^{+\infty}f\)。同时对称地取两端极限，是另一种取法，不能替代两侧分别收敛。

例如，\(\int_0^{+\infty}e^{-x}\,dx=\lim_{R\to\infty}(1-e^{-R})=1\)。虽然区间无限长，总累积仍有限。对幂函数，直接计算有限上限积分得到
\begin{equation*}
 \int_1^{+\infty}\frac{dx}{x^p}
 \quad\text{在且仅在 }p>1\text{ 时收敛，此时值为 }\frac1{p-1}.
\end{equation*}
\(p=1\)时有限积分为\(\ln R\)，其余情形由幂函数的端部行为判断。

\begin{example}
讨论\(\int_0^{+\infty}xe^{-x}\,dx\)的敛散性并求值。再考察\(f(x)=x\)：对称截取区间的积分极限，能否说明其在全实轴上的反常积分收敛？
\end{example}

\subsection{无界函数的积分}

若\(f\)在\(a\)右侧无界，但在每个\([a+\varepsilon,b]\)上可积，则定义
\begin{equation*}
 \int_a^bf(x)\,dx=\lim_{\varepsilon\to0^+}\int_{a+\varepsilon}^bf(x)\,dx,
\end{equation*}
其中右端为有限极限时称为收敛。在右端无界的情形类似；区间内部有奇点时，先在奇点两侧拆开，分别要求收敛。

对\(x^{-p}\)，有
\begin{equation*}
 \int_0^1\frac{dx}{x^p}
 \quad\text{在且仅在 }p<1\text{ 时收敛，此时值为 }\frac1{1-p}.
\end{equation*}
这与无穷远的判别方向不同：靠近零点时，指数越大增长越快；趋向无穷远时，指数越大衰减越快。

在避开零点的有限积分中令\(x=1/t\)，再取极限，就能联系这两种情形：
\begin{equation*}
 \int_0^1x^{-p}\,dx=\int_1^{+\infty}t^{p-2}\,dt
\end{equation*}
（两边按相应反常积分理解）。因此零点附近的\(p<1\)，正好对应无穷远处幂指数\(2-p>1\)。端部位置改变时，微分\(dx=-dt/t^2\)也参与改变衰减阶数。

例如，\(1/\sqrt{x}\)虽在零点附近无界，积分仍为\(2\)。另一方面，\(\int_{-1}^1dx/x\)不收敛，即使对称避开零点的两个有限积分之和始终为零。

\begin{figure}[!ht]
\centering
\includegraphics[width=.83\linewidth]{calculus-improper}
\caption{左图先截断无穷区间，再估计橙色尾部；右图先避开零点附近的奇性，再令截断位置趋近零点。}
\label{fig:calculus-improper}
\end{figure}

\begin{exercise}
\mbox{}\par\nobreak

（1）判断\(\int_0^1\ln x\,dx\)、\(\int_1^{+\infty}dx/(1+x^2)\)的敛散性，收敛时求值。

（2）分别讨论\(\int_0^1dx/(x\sqrt{x})\)与\(\int_1^{+\infty}dx/(x\sqrt{x})\)。

（3）说明为什么不能直接把\(\ln|x|\)在\(-1,1\)两端的值相减，就断定\(\int_{-1}^1dx/x=0\)。
\end{exercise}

\subsection{收敛判别}

若\(f\ge0\)，有限上限积分随上限单调不减。因此，反常积分收敛，当且仅当这些有限积分有统一上界。这使第一章的单调有界思想再次成为存在性工具。

\begin{theorem}[比较判别]\label{thm:calculus-improper-comparison}
设\(0\le f\le g\)，各函数在有限闭子区间上可积。若\(\int_a^{+\infty}g\)收敛，则\(\int_a^{+\infty}f\)收敛；若后者发散，则前者也发散。在有限奇点附近，同样可以按大小比较。
\end{theorem}

\begin{proof}
由保序性，\(\int_a^Rf\le\int_a^Rg\)。右侧有统一上界时，左侧也有；第二个结论为第一个结论的逆否命题。奇点附近同理。
\end{proof}

比较只需在足够靠近所研究端部的范围内成立，剩下的正常有限积分不影响敛散性。若正函数满足\(f/g\to c\in(0,+\infty)\)，则最终有\((c/2)g\le f\le2cg\)，所以两积分同敛散。

例如，\(x\ge1\)时\(e^{-x^2}\le e^{-x}\)，故\(\int_1^{+\infty}e^{-x^2}\,dx\)收敛，无须先求积分值。

一个积分可能在多个位置需要判别。先用正常的有限分点隔开，逐段处理；一段收敛不能抵消另一段发散。比较函数应按当前端部的行为选择。

\begin{example}
讨论\(\int_0^{+\infty}dx/[x^\alpha(1+x)]\)的敛散性。分别在零点附近与无穷远处比较幂函数；当\(\alpha=1/2\)时，用\(x=t^2\)求积分值。再讨论\(\int_0^1dx/[x^\alpha(1-x)^\beta]\)的敛散条件。
\end{example}

对可能变号的函数，可以使用Cauchy准则：\(\int_a^{+\infty}f\)收敛，当且仅当对任意\(\varepsilon>0\)，存在\(R\)，使任意\(v>u>R\)都有
\begin{equation}\label{eq:calculus-improper-cauchy}
 \left|\int_u^vf(x)\,dx\right|<\varepsilon.
\end{equation}
必要性由收敛极限直接得到。反过来，取整数上限的积分，由数列Cauchy准则得到极限，再用式\eqref{eq:calculus-improper-cauchy}比较任意足够大的上限与整数上限，便得到相同极限。这与可积性证明中先取一列、再比较任意取样的组织相似。

\subsection{绝对收敛}

\begin{definition}
若\(\int_a^{+\infty}|f(x)|\,dx\)收敛，就称\(\int_a^{+\infty}f(x)\,dx\){\heiti{}绝对收敛}。
\end{definition}

\begin{theorem}
绝对收敛的反常积分收敛。
\end{theorem}

\begin{proof}
由有限区间上的绝对值估计，\(|\int_u^vf|\le\int_u^v|f|\)。右侧由收敛性满足Cauchy条件，因此原积分也满足。
\end{proof}

例如，\(|\sin x|/x^2\le1/x^2\)，故\(\int_1^{+\infty}\sin x/x^2\,dx\)绝对收敛。原积分收敛而绝对值积分发散时，称为{\heiti{}条件收敛}。

一个典型例子是\(\int_1^{+\infty}\sin x/x\,dx\)。对\(v>u\ge1\)，分部积分得到
\begin{equation*}
 \int_u^v\frac{\sin x}{x}\,dx
 =\left.-\frac{\cos x}{x}\right|_u^v-\int_u^v\frac{\cos x}{x^2}\,dx,
 \qquad
 \left|\int_u^v\frac{\sin x}{x}\,dx\right|\le\frac2u.
\end{equation*}
所以原积分收敛。但对每个正整数\(k\)，在\([k\pi+\pi/6,k\pi+5\pi/6]\)上\(|\sin x|\ge1/2\)，故这一段的绝对值积分至少为\(1/[3(k+1)]\)。将这些互不重叠的区间相加，与调和级数比较，便知绝对值积分发散。

\begin{exercise}
\mbox{}\par\nobreak

（1）判断\(\int_1^{+\infty}dx/(x^2+\sqrt{x})\)及\(\int_0^1dx/\sqrt{x+x^2}\)的敛散性，指出比较对象。

（2）证明\(\int_0^{+\infty}e^{-x}|\cos x|\,dx\)收敛，进而说明\(\int_0^{+\infty}e^{-x}\cos x\,dx\)收敛。

（3）将上述分部积分用于\(\int_1^{+\infty}\cos x/x\,dx\)，判断它是否收敛、是否绝对收敛。
\end{exercise}

\subsection{尾部估计}

积分收敛后，还可以研究只计算有限范围造成的误差。若\(|f(x)|\le g(x)\)且\(g\)非负可积，则
\begin{equation*}
 \left|\int_a^{+\infty}f-\int_a^Rf\right|
 \le\int_R^{+\infty}g(x)\,dx.
\end{equation*}
例如，\(p>1\)时\(\int_R^{+\infty}x^{-p}\,dx=R^{1-p}/(p-1)\)。对\(R\ge1\)，又有
\begin{equation*}
 \int_R^{+\infty}e^{-x^2}\,dx
 \le\frac1R\int_R^{+\infty}xe^{-x^2}\,dx
 =\frac{e^{-R^2}}{2R}.
\end{equation*}
这既证明尾部很小，也给出选择截断位置的依据。

无穷区间上不能直接把最大误差乘以总长度。比如
\begin{equation*}
 f_n(x)=\begin{cases}\dfrac1n\left(1-\dfrac xn\right),&0\le x\le n,\\0,&x>n\end{cases}
\end{equation*}
在整个半轴上与零函数的误差不超过\(1/n\)，但积分恒为\(1/2\)。误差虽小，发生误差的范围却在扩大。

因此，比较无穷累积时，可以先固定\(R\)，控制\([a,R]\)上的误差，再分别估计两函数在\([R,+\infty)\)上的尾部。有限范围与尾部承担不同的误差控制任务。

例如，要近似\(\int_0^{+\infty}e^{-x^2}\,dx\)，先选\(R\)，使\(e^{-R^2}/(2R)\)小于给定总误差的一半；再对固定的\([0,R]\)使用中点公式，使有限区间求积误差小于另一半。先截断范围、后选分段数，两种误差便可分别验证，而不要求预先知道整个反常积分的精确值。

\subsection{积分判别法}

设\(f\)在\([1,+\infty)\)上连续、非负且单调不增。每个单位区间上的矩形与曲边面积满足
\begin{equation*}
 f(n+1)\le\int_n^{n+1}f(x)\,dx\le f(n).
\end{equation*}
累加得到
\begin{equation*}
 \int_1^{N+1}f(x)\,dx\le\sum_{n=1}^Nf(n)
 \le f(1)+\int_1^Nf(x)\,dx.
\end{equation*}
由正项级数与非负反常积分的单调有界判据，\(\sum_{n=1}^{\infty}f(n)\)与\(\int_1^{+\infty}f(x)\,dx\)同敛散。这称为{\heiti{}积分判别法}。

收敛时，还得到尾和的两侧估计：
\begin{equation*}
 \int_{N+1}^{+\infty}f(x)\,dx
 \le\sum_{n=N+1}^{\infty}f(n)
 \le\int_N^{+\infty}f(x)\,dx.
\end{equation*}
例如，\(p>1\)时，\(\sum_{n>N}n^{-p}\)介于\((N+1)^{1-p}/(p-1)\)与\(N^{1-p}/(p-1)\)之间。这接回了第一章的\(p\)级数，并进一步给出截断误差。

令\(D_N=\sum_{n=1}^Nf(n)-\int_1^Nf(x)\,dx\)。前面的估计给出\(0\le D_N\le f(1)\)，且
\begin{equation*}
 D_{N+1}-D_N=f(N+1)-\int_N^{N+1}f(x)\,dx\le0,
\end{equation*}
故\(D_N\)单调有界，必收敛。取\(f(x)=1/x\)，即得\(H_N-\ln N\to\gamma\)；调和级数与积分都发散，差值仍有极限，误差见式\eqref{eq:calculus-euler-constant-error}。

\begin{exercise}
\mbox{}\par\nobreak

（1）估计\(\sum_{n=N+1}^{\infty}1/n^2\)。要使尾和小于\(10^{-3}\)，可以选取怎样的\(N\)？

（2）讨论\(\sum_{n=2}^{\infty}1/[n(\ln n)^p]\)的敛散性，并说明被积函数在哪一段起单调不增。

（3）给出\(\int_1^{+\infty}\sin x/x\,dx\)截断到\(R\)的误差界。与绝对收敛情形的尾部比较方法有何不同？

（4）若要用数值积分计算\(\int_0^{+\infty}e^{-x^2}\,dx\)，设计一种将总误差分成有限区间求积误差与尾部误差的方案。

（5）计算\(\int_e^{+\infty}dx/[x(\ln x)^2]\)，并给出级数\(\sum_{n=2}^{\infty}1/[n(\ln n)^2]\)截断到\(n=N\)后的尾和上下界。要保证尾和小于\(1/10\)，可以取怎样的\(N\)？比较它与\(\sum1/n^2\)的截断速度。

（6）讨论\(\int_0^1x^{-p}\ln x\,dx\)的敛散性，收敛时求其值。若从\(\varepsilon\)开始积分，求遗漏\((0,\varepsilon)\)所造成的误差，并说明它怎样随\(\varepsilon\to0^+\)趋零。

（7）设\(p>0\)。用分部积分证明\(\int_1^{+\infty}\sin x/x^p\,dx\)收敛，并给出截断到\(R\)的误差界。再按\(p\)的范围判断绝对收敛与条件收敛；讨论绝对值时，可以在每个周期中选取\(|\sin x|\ge1/2\)的子区间。

（8）设\(f_n(x)=n^2xe^{-nx}\)，\(x\ge0\)。求逐点极限及每个\(f_n\)在半轴上的积分。说明为何不能直接交换积分与极限，并比较它与正文中图像向远处铺开的例子：本题的面积集中到了哪里？
\end{exercise}
